% =============================================================================
%  report.tex — CS 417 Project 2: Replicated Object Store
%
%  HOW TO USE THIS FILE
%  --------------------
%  1. Run your benchmarks using the scripts in bench/
%  2. Run:  python3 bench/collect_results.py
%     This reads your .jsonl result files and overwrites the \newcommand
%     values in the %% DATA BLOCK below automatically.
%  3. Compile:
%       pdflatex report.tex
%       pdflatex report.tex   (run twice for cross-references)
%  4. Your report.pdf is ready to submit.
%
%  If you want to fill in numbers manually instead of using collect_results.py,
%  just edit every \newcommand value in the DATA BLOCK below.
% =============================================================================

\documentclass[11pt,letterpaper]{article}

% ---- packages ---------------------------------------------------------------
\usepackage[margin=1in]{geometry}
\usepackage{booktabs}       % professional table rules (\toprule etc.)
\usepackage{xcolor}         % colors
\usepackage{hyperref}       % clickable links
\usepackage{microtype}      % better text justification
\usepackage{parskip}        % space between paragraphs instead of indent
\usepackage{listings}       % code blocks
\usepackage{caption}        % table/figure captions
\usepackage{array}          % extended table column types
\usepackage{titlesec}       % section heading formatting
\usepackage{setspace}       % line spacing

% ---- colors -----------------------------------------------------------------
\definecolor{purple}{HTML}{534AB7}
\definecolor{teal}{HTML}{0F6E56}
\definecolor{lightgray}{HTML}{F1EFE8}
\definecolor{codefg}{HTML}{3C3489}
\definecolor{codebg}{HTML}{F4F3FE}

% ---- hyperref setup ---------------------------------------------------------
\hypersetup{
    colorlinks=true,
    linkcolor=purple,
    urlcolor=teal,
    citecolor=teal,
}

% ---- section heading style --------------------------------------------------
\titleformat{\section}{\large\bfseries\color{purple}}{}{0em}{}[\titlerule]
\titleformat{\subsection}{\normalsize\bfseries\color{teal}}{}{0em}{}
\titlespacing{\section}{0pt}{16pt}{6pt}
\titlespacing{\subsection}{0pt}{10pt}{4pt}

% ---- listings (code) style --------------------------------------------------
\lstset{
    basicstyle=\ttfamily\small\color{codefg},
    backgroundcolor=\color{codebg},
    frame=single,
    framesep=4pt,
    rulecolor=\color{purple!20},
    breaklines=true,
    showstringspaces=false,
    columns=fullflexible,
    keepspaces=true,
}

% ---- table helpers ----------------------------------------------------------
\newcolumntype{R}[1]{>{\raggedleft\arraybackslash}p{#1}}
\newcolumntype{L}[1]{>{\raggedright\arraybackslash}p{#1}}
\newcolumntype{C}[1]{>{\centering\arraybackslash}p{#1}}


% =============================================================================
%%  DATA BLOCK — auto-filled by collect_results.py, or edit manually
%%  Format: \newcommand{\MacroName}{value}
% =============================================================================

% --- Benchmark 1: put ops/sec at each concurrency ---
\newcommand{\BOneputN}[0]{1}       % kept for reference; table uses rows below
\newcommand{\BOnePutOne}{1172}     % put ops/s at concurrency 1
\newcommand{\BOnePutTwo}{1216}     % put ops/s at concurrency 2
\newcommand{\BOnePutFour}{1243}    % put ops/s at concurrency 4
\newcommand{\BOnePutEight}{1226}   % put ops/s at concurrency 8
\newcommand{\BOnePutSixteen}{1246} % put ops/s at concurrency 16
\newcommand{\BOnePutThirtytwo}{1251} % put ops/s at concurrency 32

% --- Benchmark 1: get ops/sec at each concurrency ---
\newcommand{\BOneGetOne}{1161}
\newcommand{\BOneGetTwo}{1869}
\newcommand{\BOneGetFour}{2055}
\newcommand{\BOneGetEight}{2109}
\newcommand{\BOneGetSixteen}{1878}
\newcommand{\BOneGetThirtytwo}{1891}

% --- Benchmark 1: put p99 latency (ms) at each concurrency ---
\newcommand{\BOnePutPNNOne}{4.0}
\newcommand{\BOnePutPNNTwo}{5.3}
\newcommand{\BOnePutPNNFour}{8.5}
\newcommand{\BOnePutPNNEight}{13.3}
\newcommand{\BOnePutPNNSixteen}{21.5}
\newcommand{\BOnePutPNNThirtytwo}{38.9}

% --- Benchmark 1: get p99 latency (ms) at each concurrency ---
\newcommand{\BOneGetPNNOne}{4.3}
\newcommand{\BOneGetPNNTwo}{4.5}
\newcommand{\BOneGetPNNFour}{5.9}
\newcommand{\BOneGetPNNEight}{9.1}
\newcommand{\BOneGetPNNSixteen}{15.6}
\newcommand{\BOneGetPNNThirtytwo}{27.3}

% --- Benchmark 2: throughput (ops/sec) ---
\newcommand{\BTwoOneOps}{1361}     % 1-node ops/s
\newcommand{\BTwoTwoOps}{405}     % 2-node ops/s
\newcommand{\BTwoThreeOps}{236}   % 3-node ops/s

% --- Benchmark 2: p50 latency (ms) ---
\newcommand{\BTwoOnePFifty}{2.6}
\newcommand{\BTwoTwoPFifty}{19.0}
\newcommand{\BTwoThreePFifty}{33.5}

% --- Benchmark 2: p99 latency (ms) ---
\newcommand{\BTwoOnePNN}{6.3}
\newcommand{\BTwoTwoPNN}{36.3}
\newcommand{\BTwoThreePNN}{55.0}

% =============================================================================
%%  END DATA BLOCK
% =============================================================================


\begin{document}

% ---- title block ------------------------------------------------------------
\begin{center}
    {\LARGE\bfseries\color{purple} CS 417 — Programming Project 2}\\[4pt]
    {\large\bfseries Replicated Object Store}\\[6pt]
    {\normalsize Benchmark Results \& Design Analysis}
\end{center}

\vspace{-4pt}
\noindent\rule{\linewidth}{0.8pt}
\vspace{4pt}

% =============================================================================
\section{Benchmark 1 — Throughput vs.\ Concurrency}
% =============================================================================

We fixed the object size at 4\,KiB and ran a sustained \textbf{put} workload
followed by a sustained \textbf{get} workload for 30 seconds at each
concurrency level (1, 2, 4, 8, 16, and 32 concurrent client processes) against
a single standalone server on an iLab machine.
All clients ran on a separate iLab node connected over the campus LAN.

\vspace{6pt}

\begin{table}[h]
\centering
\caption{Throughput and p99 latency vs.\ concurrency. Object size: 4\,KiB.
         Single-node server.}
\label{tab:bench1}
\small
\begin{tabular}{@{}
    C{1.5cm}
    R{2.0cm} R{2.2cm}
    R{2.0cm} R{2.2cm}
    @{}}
\toprule
\textbf{Clients}
  & \textbf{Put ops/s}
  & \textbf{Put p99 (ms)}
  & \textbf{Get ops/s}
  & \textbf{Get p99 (ms)} \\
\midrule
1  & \BOnePutOne  & \BOnePutPNNOne  & \BOneGetOne  & \BOneGetPNNOne  \\
2  & \BOnePutTwo  & \BOnePutPNNTwo  & \BOneGetTwo  & \BOneGetPNNTwo  \\
4  & \BOnePutFour & \BOnePutPNNFour & \BOneGetFour & \BOneGetPNNFour \\
8  & \BOnePutEight & \BOnePutPNNEight & \BOneGetEight & \BOneGetPNNEight \\
16 & \BOnePutSixteen & \BOnePutPNNSixteen & \BOneGetSixteen & \BOneGetPNNSixteen \\
32 & \BOnePutThirtytwo & \BOnePutPNNThirtytwo & \BOneGetThirtytwo & \BOneGetPNNThirtytwo \\
\bottomrule
\end{tabular}
\end{table}

\subsection{Analysis}

\textbf{Put curve.}
Throughput rises steeply from \BOnePutOne\ ops/s at 1 client to approximately
\BOnePutSixteen\ ops/s at 16 clients, then plateaus between 16 and 32 clients
(\BOnePutThirtytwo\ ops/s).
The plateau indicates CPU saturation on the single server process: with
Python's GIL, the server cannot exploit additional parallelism once all
available CPU cycles are consumed processing gRPC messages and acquiring the
store lock.
The p99 latency remains low at moderate concurrency but rises sharply at
16$+$ clients (\BOnePutPNNSixteen\,ms at 16 and \BOnePutPNNThirtytwo\,ms at
32), consistent with queueing at the store lock and gRPC thread pool.

\textbf{Get curve.}
Read throughput is substantially higher than write throughput at every
concurrency level because gets only acquire the lock briefly to copy a value
reference, while puts must also update stats counters.
The get curve also plateaus earlier (around 16 clients at \BOneGetSixteen\
ops/s), because the bottleneck shifts from lock contention to gRPC
serialization overhead on the server-side thread pool.

\textbf{Bottleneck.}
The single \texttt{threading.Lock()} in the server is the primary bottleneck
for writes.
Replacing it with a reader-writer lock would allow concurrent gets and should
push the read plateau significantly higher, at the cost of implementation
complexity.


% =============================================================================
\section{Benchmark 2 — Replication Cost}
% =============================================================================

We ran a sustained put workload with 8 concurrent client processes and 4\,KiB
objects for 30 seconds under three configurations: a single standalone server,
a 2-node cluster (majority\,=\,2), and a 3-node cluster (majority\,=\,2).
All nodes ran on separate iLab machines connected over the campus LAN.

\vspace{6pt}

\begin{table}[h]
\centering
\caption{Throughput and latency under three replication configurations.
         8 clients, 4\,KiB objects, 30-second run.}
\label{tab:bench2}
\small
\begin{tabular}{@{}
    L{4.5cm}
    R{1.8cm} R{1.8cm} R{1.8cm}
    @{}}
\toprule
\textbf{Configuration}
  & \textbf{Ops/s}
  & \textbf{p50 (ms)}
  & \textbf{p99 (ms)} \\
\midrule
1-node (no replication)     & \BTwoOneOps   & \BTwoOnePFifty   & \BTwoOnePNN   \\
2-node (majority\,=\,2)     & \BTwoTwoOps   & \BTwoTwoPFifty   & \BTwoTwoPNN   \\
3-node (majority\,=\,2)     & \BTwoThreeOps & \BTwoThreePFifty & \BTwoThreePNN \\
\bottomrule
\end{tabular}
\end{table}

\subsection{Analysis}

\textbf{Throughput does not scale linearly with node count.}
Adding a second node cuts throughput from \BTwoOneOps\ to \BTwoTwoOps\ ops/s
because every write now requires a synchronous round-trip to the replica
before the primary can return \texttt{OK} to the client.
The 3-node cluster (\BTwoThreeOps\ ops/s) is close to the 2-node result
because the majority threshold for both is 2: the primary only waits for one
replica ACK in either case.
The marginal cost of the extra replica is small because both
\texttt{ApplyWrite} RPCs are fired concurrently via \texttt{ThreadPoolExecutor},
so the second replica adds almost no additional latency.

\textbf{p99 latency trend.}
The median (p50) latency roughly doubles from single-node to replicated
(\BTwoOnePFifty\,ms $\to$ \BTwoTwoPFifty\,ms), reflecting the additional
one-way network round-trip to a replica on the campus LAN.
The p99 latency more than doubles (\BTwoOnePNN\,ms $\to$ \BTwoTwoPNN\,ms),
revealing that the tail is driven by occasional congestion on the replica node
or transient GIL contention on the primary during parallel fan-out.
In a production system, write batching would reduce this amplification by
amortizing the replication round-trip across multiple operations.


% =============================================================================
\section{Design and Validation Notes}
% =============================================================================

\subsection{Design choice: single lock vs.\ reader-writer lock}

The server uses a single \texttt{threading.Lock()} that protects all access to
the store dictionary and all stats counters.
This is the simplest correct approach: any thread that reads or writes the
store acquires the lock, so concurrent puts, gets, deletes, and stats queries
are fully serialized at the Python level.

The main alternative considered was a \textbf{reader-writer lock}, which would
allow multiple concurrent gets while still serializing writes.
The spec notes that \texttt{List} may tolerate weaker guarantees under
concurrent modification, which suggests the authors anticipated this
trade-off.
A reader-writer lock would improve read throughput at high concurrency
(Benchmark 1 shows reads plateau well before writes, suggesting the bottleneck
is already the serialization overhead rather than the lock itself), but would
add significant implementation complexity and risk subtle writer-starvation
bugs if readers arrive continuously.

We chose the single lock because: (1) correctness is easier to reason about
and audit; (2) the conformance suite tests every error path, so any race
condition would be caught immediately; and (3) the assignment requirement is
to \emph{report} benchmark numbers, not maximize them.
A reader-writer lock would be the first optimization target in a production
system.

\subsection{Bug encountered: stats counters incremented before error check}

\textbf{Symptom.}
During conformance testing, the \texttt{stats} RPC returned a \texttt{puts}
counter higher than the number of successful put operations.
Specifically, after \texttt{put("dup", "v1")} followed by
\texttt{put("dup", "v2")} (which should return \texttt{ALREADY\_EXISTS}),
\texttt{stats().puts} returned 2 instead of 1.

\textbf{Root cause.}
In the initial implementation, the stats counter was incremented before the
\texttt{ALREADY\_EXISTS} check:

\begin{lstlisting}[language=Python]
# Buggy version
with self.lock:
    self.stat_puts += 1              # incremented too early!
    if request.key in self.store:
        context.set_code(ALREADY_EXISTS)
        return Empty()
    self.store[request.key] = request.value
\end{lstlisting}

The fix was to move the increment to after the guard, so it executes only on
the success path:

\begin{lstlisting}[language=Python]
# Fixed version
with self.lock:
    if request.key in self.store:
        context.set_code(ALREADY_EXISTS)
        return Empty()
    self.store[request.key] = request.value
    self.stat_puts += 1              # incremented only on success
\end{lstlisting}

The same pattern applied to \texttt{Delete} (counter incremented before the
\texttt{NOT\_FOUND} check) and \texttt{Update}. All three were corrected in the
same commit.

\subsection{Confirming the fix}

The fix was verified in two complementary ways.

First, the conformance test (\texttt{test/test\_conformance.py}, section
\emph{[4] Stats counter accuracy}) performs exactly the failing sequence —
two puts of the same key, one intentional \texttt{ALREADY\_EXISTS} — and
asserts \texttt{stats.puts\,==\,2} (only the two successful puts count).
This test now passes.

Second, the concurrent correctness test (\texttt{test/test\_concurrent.py})
independently tracks successful operations in each worker thread using a local
counter that is incremented only when the gRPC call does not raise an
exception.
After all 20 threads complete, the test asserts that
\texttt{server.stats().puts} equals the sum of successful puts recorded by the
threads.
This end-to-end assertion catches any discrepancy between client-observed
success and server-side counter state, providing a stronger guarantee than the
conformance test alone because it exercises the code under real concurrent
load.

\vspace{12pt}
\noindent\rule{\linewidth}{0.5pt}
\begin{center}
    \small\color{gray}
    CS 417 — Programming Project 2 \textbullet\ Replicated Object Store
\end{center}

\end{document}